\chapter{Introduction}

\section{Background}
Natural language processing (NLP) has achieved remarkable success in
high-resource languages such as English~\cite{bibid}.  However, the
majority of the world's languages are considered \emph{low-resource},
with limited digital presence and annotated datasets.

\foreignlanguage{german}{Die natürliche Sprachverarbeitung (NLP) hat
bemerkenswerte Erfolge in ressourcenreichen Sprachen wie Englisch
erzielt~\cite{bibid}.  Allerdings ist die Mehrheit der Sprachen der
Welt als \emph{ressourcenarm} einzustufen.}

\section{Research Questions}

This thesis addresses the following research questions:

\begin{enumerate}
    \item \textbf{RQ1:} How can multilingual pre-trained models be
          fine-tuned for low-resource languages with minimal labeled data?

    \item \textbf{RQ2:} What is the transfer efficiency between
          typologically distant language pairs?

    \item \textbf{RQ3:} How does the amount of target-language data
          affect transfer quality across different model architectures?

    \item \textbf{RQ4:} Can zero-shot transfer be improved by
          intermediate language selection?

    \item \textbf{RQ5:} What role does script similarity play in
          cross-lingual transfer?
\end{enumerate}

\section{Contributions}

This thesis makes the following contributions:

\begin{enumerate}
    \item A novel \textbf{adaptive transfer} method that dynamically
          selects source languages based on typological features.

    \item A comprehensive \textbf{evaluation framework} covering 15
          languages from 8 language families.

    \item An open-source \textbf{multilingual NER benchmark} with
          consistent annotation guidelines across languages.

    \item Theoretical analysis of \textbf{transfer bounds} for
          typologically diverse language pairs.

    \item Practical \textbf{recommendations} for practitioners working
          with low-resource languages.
\end{enumerate}

\section{Thesis Outline}
The remainder of this thesis is structured as follows.
\cref{chap:lit-review} reviews related work on multilingual NLP.
\cref{chap:methodology} presents our proposed approach.
Results and analysis are discussed in subsequent chapters.

\foreignlanguage{french}{La suite de cette thèse est structurée comme
suit.  Le \cref{chap:lit-review} passe en revue les travaux connexes
sur le NLP multilingue.  Le \cref{chap:methodology} présente notre
approche proposée.}

\foreignlanguage{german}{Der Rest dieser Thesis ist wie folgt
strukturiert.  \cref{chap:lit-review} behandelt verwandte Arbeiten
zum mehrsprachigen NLP.  \cref{chap:methodology} stellt unseren
vorgeschlagenen Ansatz vor.}
